\section{Experimental Setup}
\label{sec:experiments}

\subsection{Model}

We evaluate using \textbf{Gemma-2B-it}~\cite{team2024gemma}, a 2-billion parameter instruction-tuned model from Google DeepMind. We select Gemma-2B for its balance between capability and computational tractability, enabling full attention weight extraction and gradient computation on consumer hardware. The model is loaded in float32 precision with \texttt{eager} attention implementation to enable \texttt{output\_attentions=True}.

\subsection{Datasets}

\textbf{HaluEval}~\cite{li2023halueval}. A dedicated hallucination evaluation benchmark containing 10,000 question-answer pairs with human-annotated hallucination labels. Each sample includes a question, a knowledge context, a model-generated answer, and a binary hallucination tag (\texttt{yes}/\texttt{no}). We use the \texttt{qa\_samples} split.

\textbf{TriviaQA}~\cite{joshi2017triviaqa}. An open-domain question answering dataset with 95K question-answer pairs. We use the \texttt{rc.nocontext} configuration to evaluate the model's ability to answer from parametric knowledge alone, without retrieval context.

For the test evaluation, we sample $N=5$ questions per dataset with a fixed random seed ($\text{seed}=42$) for reproducibility.

\subsection{Baselines}

We compare FHI against two established hallucination detection baselines:

\textbf{Log-Probability Thresholding.} The mean log-probability of generated answer tokens is compared against a threshold $\theta = -0.5$. Answers below this threshold are classified as hallucinations~\cite{kadavath2022language}.

\textbf{Self-Consistency~\cite{wang2022self}.} The same question is posed $n=3$ times. We measure exact-match agreement among the generated answers. The hallucination risk score is $1 - \text{agreement\_ratio}$.

\subsection{Attribution Configuration}

\begin{itemize}
    \item \textbf{Primary method:} Attention rollout across all 18 transformer layers
    \item \textbf{Top-$K$:} $K = 10$ tokens selected as the attribution set
    \item \textbf{Perturbation strategies:} mask (replace with \texttt{[MASK]}), delete (remove tokens), replace (substitute with random vocabulary tokens)
\end{itemize}

\subsection{FHI Parameters}

\begin{itemize}
    \item \textbf{Weights:} $w_1 = 0.25$ (AAS), $w_2 = 0.35$ (CIS), $w_3 = 0.25$ (ESS), $w_4 = 0.15$ (HCG)
    \item \textbf{Hallucination threshold:} $\tau = 0.5$
    \item \textbf{Stability runs:} $N = 2$ (for computational efficiency on CPU)
    \item \textbf{Semantic model:} all-MiniLM-L6-v2 for embedding-based comparisons
    \item \textbf{Generation:} temperature $= 0.7$, top-$p = 0.9$, max tokens $= 512$
\end{itemize}

\subsection{Evaluation Metrics}

We evaluate hallucination detection performance using:
\begin{itemize}
    \item \textbf{Precision, Recall, F1-score} (per-class and weighted average)
    \item \textbf{AUC-ROC} where the hallucination probability is estimated as $1 - \text{FHI}$
    \item \textbf{Classification accuracy} at threshold $\tau$
\end{itemize}

\subsection{Infrastructure}

All experiments were conducted on an Apple MacBook Air (M-series, 8GB RAM) using CPU inference. Experiment tracking was performed with MLflow. The pipeline is implemented in Python 3.11 using PyTorch, HuggingFace Transformers, and Captum.
